% ============================================================
%  CONCLUSION — conclusion.tex
% ============================================================
\section{Conclusion et perspectives}
\label{sec:conclusion}

Ce rapport a présenté la conception et l'implémentation du projet Environnement,
un simulateur d'écosystème dynamique développé en Java~8 dans le cadre du cours
de Génie Logiciel et Conception. L'application constitue un système complet et
fonctionnel, capable de simuler l'évolution de biomes soumis à des événements
météorologiques variés et gouvernés par des règles de transformation écologique.

\subsection{Bilan du travail réalisé}
\label{subsec:bilan}

L'ensemble des fonctionnalités définies dans la section~\ref{sec:specification}
a été implémenté. Les sept types de biomes sont modélisés avec leurs propriétés
initiales et leur comportement propre. Les treize événements météorologiques,
statiques et mobiles, sont générés selon des formations spatiales distinctes,
déplacés et appliqués conformément aux spécifications. Les six règles de
transformation écologique évaluent correctement les propriétés des biomes à chaque
round et déclenchent les transformations attendues, telles que définies dans le
tableau~\ref{tab:transformations}. Enfin, l'interface graphique permet de
visualiser la simulation en temps réel et d'en contrôler le déroulement de manière
interactive.

Sur le plan de la conception logicielle, l'architecture modulaire à quatre modules
(\texttt{moteur}, \texttt{gui}, \texttt{config}, \texttt{util}) a démontré son
efficacité~: elle a permis un développement en parallèle, facilité les tests
d'intégration et rendu les corrections d'anomalies plus rapides à localiser. Les
patrons de conception mis en œuvre (Fabrique, Visiteur, programmation par
interfaces) ont contribué à la maintenabilité et à l'extensibilité du code, deux
critères fondamentaux identifiés dans la section~\ref{subsec:enjeux}. L'intégration
de JFreeChart a apporté une valeur analytique concrète à l'application, en
remplaçant des représentations visuelles imprécises par des graphiques rigoureux.

Le travail en équipe de trois développeurs, organisé autour du dépôt Git et de
réunions de synchronisation hebdomadaires, s'est déroulé de manière productive.
La répartition initiale des responsabilités par module, décrite dans la
section~\ref{subsec:organisation_travail}, a limité les conflits de fusion tout
en favorisant la cohésion globale du projet.

\subsection{Perspectives d'amélioration}
\label{subsec:perspectives}

Plusieurs axes d'amélioration ont été identifiés à l'issue du développement.

Le premier axe concerne la \textbf{persistance des données}. L'implémentation d'un
mécanisme de sérialisation au format JSON ou XML permettrait de sauvegarder et de
recharger l'état complet d'une simulation. Cette fonctionnalité rendrait
l'application utilisable dans un contexte pédagogique où les séances sont
fragmentées.

Le deuxième axe porte sur l'\textbf{extensibilité dynamique}. Permettre le
chargement de nouveaux types de biomes ou de nouvelles règles de transformation à
partir de fichiers de configuration externes, sans recompilation, constituerait
une amélioration architecturale significative. Un mécanisme de plugin fondé sur
l'API \textit{Service Loader} de Java serait une solution
adaptée à cette évolution.

Le troisième axe vise l'\textbf{amélioration de l'interface graphique}. Une
refonte vers une mise en page fluide et redimensionnable, recourant par exemple à
la bibliothèque JavaFX, permettrait de lever la contrainte de résolution fixe
identifiée dans la section~\ref{subsec:notions_base}. Cette refonte améliorerait
également l'expérience utilisateur sur les écrans de petite taille.

Enfin, un quatrième axe envisage l'\textbf{enrichissement du modèle de simulation}.
L'introduction de chaînes alimentaires entre espèces animales virtuelles peuplant
les biomes, ou la modélisation de phénomènes climatiques globaux (cycles
hydrologiques, effet de serre simplifié), accroîtrait considérablement le réalisme
et la valeur pédagogique du simulateur.

Ces perspectives confirment que le projet Environnement, dans son état actuel,
constitue une base solide et extensible sur laquelle des développements futurs
peuvent s'appuyer.
